% Imporditud: tools/import_eksperimendid.py
% Allikas: Eesti füüsikaolümpiaadi eksperimendi ülesannete
%   kogu (Rael Kalda, 2023), ülesanne Ü149, lahendus L149.
\setAuthor{Erkki Tempel}
\setRound{lõppvoor}
\setYear{2017}
\setNumber{G E1}
\setDifficulty{4}
\setTopic{vedelike mehaanika}
\setVahendid{erinevast materjalist kehad A ja B, niit, purk veega, puidust mõõ-}
\setPunktid{12}
\setAllikas{Eksperimendi kogu Ü149, lahendus lk 111}
\setLahendusseis{toores}

\prob{Ruumalad}
Määrake võimalikult täpselt erinevast materjalist kehade A ja B tiheduste suhe $\rho$A . $\rho$B

Materjalide tihedus on suurem kui vee tihedus.

tejoonlaud.

\hint

\solu
Kasutame joonlauda kangina. Toetuspunktiks on lauaserv. Leiame joonlaua massikeskme. Edaspidi on joonlaua massikese täpselt laua serval. Leiame kehade jõuõlgade lA1 ja lB1 pikkused toetuspunktist. Siis kangireegli järgi kehtib

mAlA1 = mB lB1

Nüüd uputame keha A keha vette ning leiame uuesti jõuõlgade lA2 ja lB2 kaugused toetuspunktist. Kehale A mõjub ka üleslükkejõud Fy = $\rho$vesigVA ja kangireeglist saame

(mA $- \rho$vesiVA)lA2 = mBlB2

Nüüd võtame keha A veest välja ja uputame keha B keha vette ning leiame uuesti jõuõlgade lA3 ja lB3 kaugused toetuspunktist. Kehale B mõjub üleslükkejõud Fy = $\rho$vesigVB ja kangireeglist saame

mAlA3 = (mB $- \rho$vesiVB)lB3

Tiheduste suhe on ruumalade ja masside suhe

$\rho$A = mAVB $\rho$B mBVA

Kõikidest eelnevatest võrrandistest koosnevast süsteemist tiheduste suhte avalda-

des saame lA3 lB1

$\rho$A = 1$-$ . lA1 lB3 $\rho$B 1$-$ lA1 lB2 lA2 lB1
\probend
